\begin{ThreePartTable}

    \renewcommand\TPTminimum{\textwidth}

    \setlength\LTleft{0pt}
    \setlength\LTright{0pt}
    \setlength\tabcolsep{3pt}

    \begin{TableNotes}
        \item \textbf{Nota.} Adaptado de: \textcite[p. 44]{AENOR:2016} [traducido del inglés].
    \end{TableNotes}

    \begin{longtable}{@{}P{0.2\textwidth}P{0.2\textwidth}P{0.6\textwidth} @{}}

        \caption{Utilidad de la información de mantenimiento} \label{tab:maintenance_information_utility} \\

        \toprule
        \textbf{Información a ser recogida} &
            \textbf{Prioridad en la recogida de datos} &
                \textbf{Ejemplos} \\
        \toprule
        \endfirsthead
        
        \multicolumn{3}{r}{\textit{(Continuación de la tabla~\ref{tab:maintenance_information_utility})}} \\
        \toprule
        \textbf{Información a ser recogida} &
            \textbf{Prioridad en la recogida de datos} &
                \textbf{Ejemplos} \\
        \toprule
        \endhead

        \midrule[\heavyrulewidth]
        \multicolumn{3}{r}{\textit{Continúa en la siguiente página}}\\
        \endfoot

        \bottomrule
        \insertTableNotes\\
        \endlastfoot

        Mantenimiento correctivo &
            Requerida &
                \begin{itemize}[left=1em, nosep]
                    \item Tiempo de reparación (MTTR o MRT)
                    \item Cantidad de mantenimiento correctivo
                    \item Estrategia de reemplazo o reparación
                \end{itemize} \\

        \midrule

        Mantenimiento preventivo &
            Opcional &
                \begin{itemize}[left=1em, nosep]
                    \item Historial de la vida útil del equipo
                    \item Cantidad total de recursos usados
                    \item Tiempo muerto total
                    \item Efecto del mantenimiento preventivo en la tasa de falla
                    \item Balance entre mantenimiento preventivo y correctivo
                \end{itemize} \\

        \midrule

        Mantenimiento programado &
            Opcional &
                \begin{itemize}[left=1em, nosep]
                    \item Diferencia entre el mantenimiento estimado y ejecutado (backlog)
                    \item Actualización de los programas de mantenimiento basado en la experiencia (métodos, recursos e intervalos)
                \end{itemize} \\

    \end{longtable}

\end{ThreePartTable}
\fixvspace{0}